% ASSERT_CONVENTION: natural_units=natural, metric_signature=mostly_minus, fourier_convention=physics, coupling_convention=alpha_s, generator_normalization=T(fund)=1/2, covariant_derivative_sign=D_mu=partial_mu-igA
%
% Exercise Set 6: Scale Matching -- v and Lambda_QCD from Dual Parameters
% Part of: The Dualised Standard Model -- Part III Exercises
%
% Conventions (Seiberg hep-th/9411149, unchanged from Lectures 1-8):
%   T(fund) = 1/2 per Weyl fermion
%   b_0 = (11N_c - 2N_f)/3 for N_f Dirac flavours
%   alpha_s = g_s^2 / (4 pi)
%   Epistemic: all non-SUSY duality claims carry conjectural qualifiers (P1)
%
\documentclass[11pt,a4paper]{article}
% ASSERT_CONVENTION: natural_units=natural, metric_signature=mostly_minus, fourier_convention=physics, coupling_convention=alpha_s, generator_normalization=T(fund)=1/2, covariant_derivative_sign=D_mu=partial_mu-igA
%
% CONVENTIONS (Seiberg hep-th/9411149)
% Metric: (+,-,-,-) (mostly minus)
% T(fund) = 1/2 per Weyl fermion
% b_0 = (11N_c - 2N_f)/3 for N_f Dirac flavours
% D_mu = partial_mu - ig A_mu^a T^a
% A(fund) = 1 (cubic anomaly coefficient per fundamental Weyl fermion)
% alpha_s = g_s^2 / (4 pi)
% R[Q] = (N_f - N_c)/N_f; R[W] = 2
% N_f counts Dirac flavour pairs (each = 2 Weyl fermions)
%
% This file is \input{} by all lecture files.

%% ---------- Document class ----------
% (loaded by the wrapper; preamble only provides packages and macros)

%% ---------- Standard packages ----------
\usepackage{amsmath,amssymb,amsthm,mathtools}
\usepackage{hyperref}
\usepackage[capitalise,noabbrev]{cleveref}
\usepackage{enumitem}
\usepackage{booktabs}
\usepackage{array}
\usepackage{xcolor}
\usepackage{tikz}
\usetikzlibrary{arrows.meta,decorations.markings,positioning,calc}
\usepackage{fancybox}
\usepackage{tcolorbox}
\tcbuselibrary{skins,breakable}
\usepackage{slashed}              % Feynman slash notation
\usepackage{cite}                 % compressed citation lists

%% ---------- Hyperref configuration ----------
\hypersetup{
  colorlinks=true,
  linkcolor=blue!70!black,
  citecolor=green!50!black,
  urlcolor=blue!80!black,
  pdfauthor={Part III Lecture Notes},
  pdftitle={The Dualised Standard Model}
}

%% ---------- Theorem environments ----------
\theoremstyle{plain}
\newtheorem{theorem}{Theorem}[section]
\newtheorem{lemma}[theorem]{Lemma}
\newtheorem{proposition}[theorem]{Proposition}
\newtheorem{corollary}[theorem]{Corollary}

\theoremstyle{definition}
\newtheorem{definition}[theorem]{Definition}
\newtheorem{example}[theorem]{Example}
\newtheorem{exercise}{Exercise}[section]

\theoremstyle{remark}
\newtheorem{remark}[theorem]{Remark}

%% ---------- Physics macros: gauge groups and representations ----------
\newcommand{\Nc}{N_{\!c}}
\newcommand{\Nf}{N_{\!f}}
\newcommand{\SU}[1]{\mathrm{SU}(#1)}
\newcommand{\UU}[1]{\mathrm{U}(#1)}

% Representations
\newcommand{\fund}{\square}
\newcommand{\antifund}{\overline{\square}}
\newcommand{\adj}{\mathrm{adj}}
\newcommand{\antisym}{\wedge^{\!2}}        % antisymmetric rank-2
\newcommand{\sym}{\mathrm{S}_2}           % symmetric rank-2

%% ---------- Physics macros: group theory constants ----------
% Dynkin index: Tr(T^a T^b) = T(R) delta^{ab}
\newcommand{\Tfund}{\tfrac{1}{2}}          % T(fund) = 1/2 per Weyl fermion
\newcommand{\Tadj}{\Nc}                    % T(adj) = N_c
\newcommand{\Cfund}{\frac{\Nc^2 - 1}{2\Nc}}  % C_2(fund)
\newcommand{\Cadj}{\Nc}                    % C_2(adj)

% Cubic anomaly coefficient: A(R) = Tr_R[T^a {T^b, T^c}] with A(fund) = 1
\newcommand{\Afund}{1}

%% ---------- Physics macros: beta function ----------
\newcommand{\bzero}{b_0}
\newcommand{\bone}{b_1}
\newcommand{\alphas}{\alpha_s}
\newcommand{\gs}{g_s}

%% ---------- Physics macros: SUSY / R-charge ----------
\newcommand{\Rcharge}[1]{R[#1]}
\newcommand{\Wsup}{W}                      % superpotential

%% ---------- Physics macros: covariant derivative and field strength ----------
\newcommand{\Dmu}{D_\mu}
\newcommand{\Fmn}{F_{\mu\nu}}
\newcommand{\Ftilde}{\widetilde{F}_{\mu\nu}}

%% ---------- Physics macros: common abbreviations ----------
\newcommand{\ie}{\textit{i.e.}}
\newcommand{\eg}{\textit{e.g.}}
\newcommand{\cf}{\textit{cf.}}
\newcommand{\IR}{\ensuremath{\mathrm{IR}}}
\newcommand{\UV}{\ensuremath{\mathrm{UV}}}
\newcommand{\AF}{\ensuremath{\mathrm{AF}}}
\newcommand{\BZ}{\ensuremath{\mathrm{BZ}}}
\newcommand{\NSVZ}{\ensuremath{\mathrm{NSVZ}}}
\newcommand{\SQCD}{\ensuremath{\mathrm{SQCD}}}
\newcommand{\QCD}{\ensuremath{\mathrm{QCD}}}
\newcommand{\SM}{\ensuremath{\mathrm{SM}}}
\newcommand{\Tr}{\mathrm{Tr}}

%% ---------- Convention box macro ----------
\newtcolorbox{conventionbox}[1][]{%
  enhanced,
  colback=blue!5!white,
  colframe=blue!60!black,
  fonttitle=\bfseries,
  title={Conventions},
  #1
}

%% ---------- Comparison table style ----------
\newtcolorbox{comparisontable}[1][]{%
  enhanced,
  colback=orange!5!white,
  colframe=orange!70!black,
  fonttitle=\bfseries,
  title={Comparison},
  #1
}

%% ---------- Warning / important box ----------
\newtcolorbox{warningbox}[1][]{%
  enhanced,
  colback=red!5!white,
  colframe=red!70!black,
  fonttitle=\bfseries,
  title={Important},
  #1
}

%% ---------- Preview / forward-reference box ----------
\newtcolorbox{previewbox}[1][]{%
  enhanced,
  colback=green!5!white,
  colframe=green!50!black,
  fonttitle=\bfseries,
  title={Preview},
  #1
}

%% ---------- Page layout ----------
\usepackage[margin=2.5cm]{geometry}
\setlength{\parskip}{0.4em}
\setlength{\parindent}{0em}

%% ---------- Section formatting ----------
\usepackage{titlesec}
\titleformat{\section}{\Large\bfseries}{\thesection}{1em}{}
\titleformat{\subsection}{\large\bfseries}{\thesubsection}{1em}{}

%% ---------- End of preamble ----------


\title{\textbf{Exercise Set 6: Scale Matching --- $v$ and $\Lambda_{\QCD}$
  from Dual Parameters}}
\author{The Dualised Standard Model --- Part III Exercises}
\date{Lent Term 2027}

\begin{document}
\maketitle

\begin{conventionbox}[title={Conventions for these exercises}]
\renewcommand{\arraystretch}{1.3}
\begin{tabular}{@{}l l@{}}
\toprule
\textbf{Quantity} & \textbf{Convention} \\
\midrule
Units & $\hbar = c = 1$ (natural) \\
Metric & $\eta_{\mu\nu} = \mathrm{diag}(+1,-1,-1,-1)$ (mostly minus) \\
Beta function & $\bzero = (11\Nc - 2\Nf)/3$ for $\Nf$ Dirac flavours \\
Coupling & $\alphas = g_s^2/(4\pi)$ \\
Planck mass & $M_{\text{Pl}} = 1.22 \times 10^{19}\;\text{GeV}$ \\
Epistemic & ``if the conjectured duality holds'' for all non-SUSY claims (P1) \\
\bottomrule
\end{tabular}

\medskip

\noindent References: Lecture~1 (\S\S2--3, RG running);
Lecture~8 (\S\S4--5, scale matching);
Cacciapaglia--Sannino~\cite{CacciapagliaSannino2024}.

\medskip

\noindent\textbf{Epistemic note.}
All results in this exercise assume the conjectured non-SUSY
duality holds.  The one-loop QCD running is standard textbook
physics; the duality interpretation of the input coupling
$\alphas(\Lambda_D)$ is conjectural.
\end{conventionbox}

\bigskip


%% ========================================================================
%%  PROBLEM 1: Scale Matching -- v ~ 246 GeV and Lambda_QCD ~ 200 MeV
%% ========================================================================
\section*{Problem 1: Scale Matching in the Dualised Standard Model
  \hfill $\star\star\star$ \normalsize(25 min, 25 marks)}
\addcontentsline{toc}{section}{Problem 1: Scale matching -- $v$ and
  $\Lambda_{\QCD}$}
\label{prob:scale-matching}

\noindent\textbf{Background.}
In the dualised Standard Model, the \textbf{duality scale}
$\Lambda_D$ is the energy at which the electric $\SU{2}$ gauge
theory confines and the magnetic (SM) description becomes weakly
coupled.  Below $\Lambda_D$, the SM gauge couplings run according
to the usual perturbative renormalisation group equations.  The
duality scale is a \textbf{free parameter} of the framework; the
estimate $\Lambda_D \sim 10^{11}\;\text{GeV}$ is a parametric
consistency check (Lecture~8, \S4.1).

This exercise derives two physical scales from the duality scale:
\begin{enumerate}
  \item The electroweak (Fermi) scale $v \approx 246\;\text{GeV}$
    from the seesaw-like relation
    $v \sim \Lambda_D^2 / M_{\text{Pl}}$.
  \item The QCD confinement scale
    $\Lambda_{\QCD} \approx 200\;\text{MeV}$
    from one-loop RG running of $\alphas$ below $\Lambda_D$.
\end{enumerate}

\medskip

\noindent\textbf{Questions.}
\begin{enumerate}[label=(\alph*)]

  \item \textbf{[5 marks]}
    The one-loop beta function for $\SU{3}_c$ QCD is:
    \begin{equation}\label{eq:beta-coeff}
      \bzero \;=\; \frac{11\Nc - 2\Nf}{3}\,.
    \end{equation}
    \begin{enumerate}[label=(\roman*)]
      \item With $\Nc = 3$ colours and $\Nf = 6$ quark flavours active
        above the top mass threshold, compute $\bzero$.

      \item The one-loop running of the strong coupling is
        (Lecture~1, Eq.~(3)):
        \begin{equation}\label{eq:one-loop-running}
          \alphas(\mu)
          \;=\; \frac{\alphas(\mu_0)}{%
            1 + \dfrac{\bzero}{2\pi}\,
            \alphas(\mu_0)\,
            \ln\!\bigl(\mu / \mu_0\bigr)}\,.
        \end{equation}
        The QCD confinement scale $\Lambda_{\QCD}$ is defined by the
        condition $\alphas(\Lambda_{\QCD}) \to \infty$, which gives:
        \begin{equation}\label{eq:lambda-qcd-def}
          \Lambda_{\QCD}
          \;=\; \mu_0 \;\exp\!\Bigl(
            -\frac{2\pi}{\bzero\,\alphas(\mu_0)}
          \Bigr)\,.
        \end{equation}
        Verify this by substituting~\eqref{eq:lambda-qcd-def}
        into~\eqref{eq:one-loop-running} and showing that the
        denominator vanishes at $\mu = \Lambda_{\QCD}$.

      \item State the dimensional analysis check:
        $[\Lambda_{\QCD}] = [\mu_0] = [\text{energy}]$.
    \end{enumerate}

  \item \textbf{[8 marks]}
    At the duality scale $\Lambda_D$, the magnetic and electric gauge
    couplings are related by the duality relation
    (Lecture~8, Eq.~(16)):
    \begin{equation}\label{eq:duality-coupling}
      \frac{1}{\alpha_e(\Lambda_D)}
      \;\sim\; \alpha_m(\Lambda_D)
      \;=\; \alphas(\Lambda_D)\,.
    \end{equation}
    Strong coupling in the electric theory
    ($\alpha_e(\Lambda_D) \gg 1$) corresponds to weak coupling in the
    magnetic theory ($\alphas(\Lambda_D) \ll 1$).

    \begin{enumerate}[label=(\roman*)]
      \item For the electric $\SU{2}$ gauge theory underlying the SM
        QCD sector, the field content includes 6 Weyl fermions in the
        fundamental representation and one Weyl fermion in the adjoint
        (from the gaugino of the progenitor SUSY theory).  Compute the
        electric one-loop coefficient
        \[
          \bzero^e
          \;=\; \frac{11 N_c^e}{3}
          \;-\; \frac{2}{3}\,n_f^e\, T(\fund)
          \;-\; \frac{2}{3}\, T(\adj)\,,
        \]
        with $N_c^e = 2$, $n_f^e = 6$ Weyl fundamentals, $T(\fund) = 1/2$, and
        $T(\adj) = N_c^e = 2$.
        Argue that the electric theory \emph{is} asymptotically
        free ($\bzero^e > 0$), which is consistent with the duality
        requirement that the electric theory becomes strongly coupled
        in the infrared (the coupling grows as $\mu$ decreases).

      \item Using the measured value $\alphas(M_Z) = 0.118$ at
        $M_Z = 91.2\;\text{GeV}$, first run the coupling up to
        $\Lambda_D = 10^{11}\;\text{GeV}$ using one-loop RG
        (Eq.~\eqref{eq:one-loop-running} with $\bzero = 7$) to
        determine $\alphas(\Lambda_D)$.  Then use
        Eq.~\eqref{eq:lambda-qcd-def} with $\mu_0 = \Lambda_D$
        to compute $\Lambda_{\QCD}$ numerically.  Show all
        intermediate steps.

      \item Compare your result with the measured value
        $\Lambda_{\QCD} \approx 200\;\text{MeV}$.  Identify two
        sources of $\mathcal{O}(1)$ uncertainty in the estimate.
    \end{enumerate}

  \item \textbf{[7 marks]}
    The electroweak scale arises from Planck-suppressed operators in the
    electric theory.  As shown in Lecture~8, \S2, the scalar
    mass-squared for the leading Higgs doublet is:
    \begin{equation}\label{eq:scalar-mass}
      \mu^2 \;\sim\; \xi\, \frac{\Lambda_D^4}{M_{\text{Pl}}^2}\,,
    \end{equation}
    where $\xi$ is an $\mathcal{O}(1)$ dimensionless coefficient from
    the Planck-scale operators.  The electroweak VEV is
    $v = \mu / \sqrt{\lambda}$, where $\lambda$ is the Higgs quartic
    coupling.

    \begin{enumerate}[label=(\roman*)]
      \item Show that, up to $\mathcal{O}(1)$ factors,
        \begin{equation}\label{eq:fermi-seesaw}
          v \;\sim\; \frac{\Lambda_D^2}{M_{\text{Pl}}}\,.
        \end{equation}
        This is a \textbf{seesaw-like relation}: the electroweak
        hierarchy $v / M_{\text{Pl}} \sim 10^{-17}$ arises from the
        intermediate hierarchy
        $\Lambda_D / M_{\text{Pl}} \sim 10^{-8}$.

      \item For $\Lambda_D = 10^{11}\;\text{GeV}$ and
        $M_{\text{Pl}} = 1.22 \times 10^{19}\;\text{GeV}$, compute
        $v$ numerically from Eq.~\eqref{eq:fermi-seesaw}.  Compare
        with $v_{\text{measured}} = 246\;\text{GeV}$.

      \item Determine the range of $\Lambda_D$ that gives
        $v$ in the range $100$--$1000\;\text{GeV}$.

      \item Verify dimensional consistency:
        $[\Lambda_D^2 / M_{\text{Pl}}] = [\text{energy}]$.
    \end{enumerate}

  \item \textbf{[5 marks]}
    \textbf{Summary and discussion.}

    \begin{enumerate}[label=(\roman*)]
      \item Complete the following input/output table for the dualised
        SM scale-matching framework:

        \renewcommand{\arraystretch}{1.3}
        \begin{center}
        \begin{tabular}{@{} l l l @{}}
        \toprule
        \textbf{Quantity} & \textbf{Role} & \textbf{Value} \\
        \midrule
        $\Lambda_D$ & \rule{3cm}{0.4pt} & $\sim 10^{11}\;\text{GeV}$ \\
        $M_{\text{Pl}}$ & \rule{3cm}{0.4pt}
          & $1.22 \times 10^{19}\;\text{GeV}$ \\
        $\alphas(\Lambda_D)$ & \rule{3cm}{0.4pt} & $\sim 0.03$ \\
        $v$ & \rule{3cm}{0.4pt} & $\sim 246\;\text{GeV}$ \\
        $\Lambda_{\QCD}$ & \rule{3cm}{0.4pt}
          & $\sim 200\;\text{MeV}$ \\
        $\Lambda_S$ & \rule{3cm}{0.4pt} & $\sim \text{few TeV}$ \\
        \bottomrule
        \end{tabular}
        \end{center}

      \item The three energy scales of the dualised SM are
        $\Lambda_S \sim \text{TeV}$ (quasi-SUSY spectrum mass),
        $\Lambda_D \sim 10^{11}\;\text{GeV}$ (duality scale), and
        optionally $\Lambda_{\text{eGUT}}$ (electric grand unification).
        Present the hierarchy
        $\Lambda_{\QCD} \ll v \ll \Lambda_S \ll \Lambda_D
        \ll M_{\text{Pl}}$
        and briefly state which scales are inputs (free parameters) and
        which are outputs (derived quantities).

      \item Is the relation $v \sim \Lambda_D^2 / M_{\text{Pl}}$ a
        \textbf{prediction} or a \textbf{parametric consistency check}?
        Justify your answer.
    \end{enumerate}

\end{enumerate}

\medskip

\noindent\textit{Hints:}
\begin{itemize}
  \item For (a)(ii): set the denominator of
    Eq.~\eqref{eq:one-loop-running} to zero and solve for
    $\mu$.
  \item For (b)(ii): first evolve $\alphas(M_Z) = 0.118$ to
    $\Lambda_D$ using Eq.~\eqref{eq:one-loop-running}, then compute
    $2\pi/(\bzero \cdot \alphas(\Lambda_D))$ and take the exponential.
  \item For (c)(iii): invert the seesaw relation to express
    $\Lambda_D$ as a function of $v$.
\end{itemize}

\bigskip
\hrule
\bigskip


%% ========================================================================
%%  SOLUTION TO PROBLEM 1
%% ========================================================================
\begin{proof}[\textbf{Solution to Problem 1}]

\textbf{(a) QCD beta function and $\Lambda_{\QCD}$ definition.}

\textit{(i) Beta function coefficient.}

With $\Nc = 3$ and $\Nf = 6$:
\[
  \bzero \;=\; \frac{11 \times 3 - 2 \times 6}{3}
  \;=\; \frac{33 - 12}{3}
  \;=\; \frac{21}{3}
  \;=\; 7\,.
\]
\[
  \boxed{\bzero = 7\;\;\text{for } \SU{3},\; \Nf = 6\,.}
  \quad\checkmark
\]

Since $\bzero > 0$, QCD is asymptotically free with 6 flavours.

\textit{(ii) Verification of $\Lambda_{\QCD}$ definition.}

From Eq.~\eqref{eq:one-loop-running}, $\alphas(\mu) \to \infty$
when the denominator vanishes:
\[
  1 + \frac{\bzero}{2\pi}\, \alphas(\mu_0)\,
  \ln\!\bigl(\mu/\mu_0\bigr) \;=\; 0\,.
\]
Solving for $\mu$:
\[
  \ln\!\bigl(\mu/\mu_0\bigr)
  \;=\; -\frac{2\pi}{\bzero\, \alphas(\mu_0)}\,,
\]
\[
  \mu \;=\; \mu_0\, \exp\!\Bigl(
    -\frac{2\pi}{\bzero\, \alphas(\mu_0)}
  \Bigr)
  \;\equiv\; \Lambda_{\QCD}\,.
  \quad\checkmark
\]

\textit{(iii) Dimensional analysis.}

$\Lambda_{\QCD} = \mu_0 \times \exp(\ldots)$.  The exponential is
dimensionless, so
$[\Lambda_{\QCD}] = [\mu_0] = [\text{energy}]$.
\quad$\checkmark$

\bigskip

\textbf{(b) $\Lambda_{\QCD}$ from the duality scale.}

\textit{(i) Electric beta function.}

For the electric $\SU{2}$ theory with $n_f^e = 6$ Weyl fundamentals
and 1 Weyl adjoint:
\begin{align}
  \bzero^e
  &\;=\; \frac{11 \times 2}{3}
  \;-\; \frac{2}{3} \times 6 \times \frac{1}{2}
  \;-\; \frac{2}{3} \times 2
  \nonumber \\
  &\;=\; \frac{22}{3} - 2 - \frac{4}{3}
  \;=\; \frac{22 - 4}{3} - 2
  \;=\; 6 - 2 \;=\; 4\,.
\end{align}

Since $\bzero^e > 0$, the electric theory is \emph{asymptotically
free}: it becomes strongly coupled in the IR (at $\Lambda_D$) and
weakly coupled in the UV.  This is precisely the requirement for
duality---the electric theory confines at $\Lambda_D$, and the
magnetic (SM) description takes over below this scale.

\begin{remark}[Consistency check]
  The electric theory confining at $\Lambda_D$ requires that the
  electric coupling $\alpha_e$ grows as the energy decreases toward
  $\Lambda_D$. With $b_0^e > 0$ (asymptotic freedom), $\alpha_e$
  increases in the IR, consistent with the duality picture.
\end{remark}

\textit{(ii) Numerical computation of $\Lambda_{\QCD}$.}

With $\alphas(\Lambda_D) = 0.3$, $\Lambda_D = 10^{11}\;\text{GeV}$,
and $\bzero = 7$:

\textbf{Step 1:} Compute the exponent:
\[
  \frac{2\pi}{\bzero\, \alphas(\Lambda_D)}
  \;=\; \frac{2\pi}{7 \times 0.3}
  \;=\; \frac{6.283}{2.1}
  \;\approx\; 2.99\,.
\]

\textbf{Step 2:} Compute $\Lambda_{\QCD}$:
\[
  \Lambda_{\QCD}
  \;=\; 10^{11}\;\text{GeV} \times
  \exp(-2.99)
  \;=\; 10^{11}\;\text{GeV} \times 0.050
  \;\approx\; 5 \times 10^{9}\;\text{GeV}\,.
\]

This is far too large!  The issue is that $\alphas(\Lambda_D) = 0.3$
is too large for the coupling at $10^{11}\;\text{GeV}$.  Let us
reconsider: the duality relation gives
$\alphas(\Lambda_D) \sim 1/\alpha_e$, and for $\alpha_e \sim
\mathcal{O}(1)$ we should use a smaller magnetic coupling.

\textbf{Re-evaluation with self-consistent approach.}
The measured value $\alphas(M_Z) = 0.118$ at $M_Z = 91.2\;\text{GeV}$
can be evolved to $\Lambda_D$ using Eq.~\eqref{eq:one-loop-running}:
\[
  \alphas(\Lambda_D)
  \;=\; \frac{0.118}{%
    1 + \dfrac{7}{2\pi} \times 0.118 \times
    \ln\!\bigl(10^{11}/91.2\bigr)}\,.
\]
Computing the logarithm:
\[
  \ln\!\bigl(10^{11}/91.2\bigr)
  \;=\; \ln(1.10 \times 10^{9})
  \;=\; 9 \ln 10 + \ln 1.10
  \;\approx\; 20.7 + 0.1
  \;=\; 20.8\,.
\]
The denominator:
\[
  1 + \frac{7}{2\pi} \times 0.118 \times 20.8
  \;=\; 1 + 1.114 \times 20.8 \times 0.118
  \;=\; 1 + \frac{7 \times 0.118 \times 20.8}{6.283}\,.
\]
Let us compute carefully:
\[
  \frac{7 \times 0.118}{2\pi}
  \;=\; \frac{0.826}{6.283}
  \;\approx\; 0.1315\,.
\]
\[
  0.1315 \times 20.8 \;\approx\; 2.74\,.
\]
\[
  \alphas(\Lambda_D)
  \;\approx\; \frac{0.118}{1 + 2.74}
  \;=\; \frac{0.118}{3.74}
  \;\approx\; 0.032\,.
\]

\textbf{Step 1 (revised):} Compute the exponent with
$\alphas(\Lambda_D) \approx 0.032$:
\[
  \frac{2\pi}{\bzero\, \alphas(\Lambda_D)}
  \;=\; \frac{6.283}{7 \times 0.032}
  \;=\; \frac{6.283}{0.224}
  \;\approx\; 28.0\,.
\]

\textbf{Step 2 (revised):} Compute $\Lambda_{\QCD}$:
\[
  \Lambda_{\QCD}
  \;=\; 10^{11}\;\text{GeV} \times
  \exp(-28.0)\,.
\]
\[
  \exp(-28.0) \;\approx\; 6.9 \times 10^{-13}\,.
\]
\[
  \Lambda_{\QCD}
  \;\approx\; 10^{11} \times 6.9 \times 10^{-13}\;\text{GeV}
  \;=\; 69\;\text{MeV}\,.
\]
\[
  \boxed{\Lambda_{\QCD}^{(\Nf=6)}
  \;\approx\; 70\;\text{MeV}\,.}
\]

This is within a \textbf{factor of 3} of the measured
$\Lambda_{\QCD} \approx 200\;\text{MeV}$---a satisfactory result
for a one-loop estimate with no threshold corrections.

\textit{(iii) Sources of $\mathcal{O}(1)$ uncertainty.}

\begin{enumerate}
  \item \textbf{Threshold corrections:} The calculation uses
    $\Nf = 6$ and $\bzero = 7$ throughout, but the number of active
    flavours decreases at each quark mass threshold (from $\Nf = 6$
    above $m_t$ to $\Nf = 5$ below $m_t$, then $\Nf = 4$ below
    $m_b$, etc.).  At each threshold, $\bzero$ increases, which
    \emph{slows} the running and increases $\Lambda_{\QCD}$.  A proper
    threshold matching would bring the estimate closer to
    $200\;\text{MeV}$.

  \item \textbf{Duality relation uncertainty:} The matching condition
    $\alphas(\Lambda_D)$ depends on the duality relation, which is an
    $\mathcal{O}(1)$ statement.  A $30\%$ change in
    $\alphas(\Lambda_D)$ shifts $\Lambda_{\QCD}$ by more than an
    order of magnitude (because of the exponential sensitivity).

  \item \textbf{Two-loop corrections:} The two-loop beta function
    shifts $\Lambda_{\QCD}$ by $\sim 10$--$20\%$ at the relevant
    energy scales.
\end{enumerate}

\bigskip

\textbf{(c) The Fermi scale from the duality scale.}

\textit{(i) Seesaw relation.}

From Eq.~\eqref{eq:scalar-mass}, the Higgs mass parameter is
$\mu^2 \sim \xi\, \Lambda_D^4 / M_{\text{Pl}}^2$.  The electroweak
VEV is $v = \mu/\sqrt{\lambda}$, where $\lambda$ is the Higgs
quartic coupling (an $\mathcal{O}(1)$ number).  Therefore:
\[
  v \;\sim\; \frac{\mu}{\sqrt{\lambda}}
  \;\sim\; \frac{1}{\sqrt{\lambda}}\,
  \sqrt{\xi}\,
  \frac{\Lambda_D^2}{M_{\text{Pl}}}\,.
\]
Absorbing the $\mathcal{O}(1)$ factors $\sqrt{\xi}/\sqrt{\lambda}$:
\[
  \boxed{v \;\sim\; \frac{\Lambda_D^2}{M_{\text{Pl}}}\,.}
\]

\textbf{Dimensional check:}
$[\Lambda_D^2/M_{\text{Pl}}] = [\text{GeV}]^2 / [\text{GeV}]
= [\text{GeV}]$.
\quad$\checkmark$

\textbf{Physical interpretation.}  This is a seesaw-like relation: the
electroweak hierarchy $v/M_{\text{Pl}} \sim 10^{-17}$ arises as the
\emph{square} of the intermediate hierarchy
$\Lambda_D/M_{\text{Pl}} \sim 10^{-8}$:
\[
  \frac{v}{M_{\text{Pl}}}
  \;\sim\; \biggl(\frac{\Lambda_D}{M_{\text{Pl}}}\biggr)^{\!2}\,.
\]

\textit{(ii) Numerical evaluation.}

\[
  v \;\sim\; \frac{(10^{11})^2}{1.22 \times 10^{19}}\;\text{GeV}
  \;=\; \frac{10^{22}}{1.22 \times 10^{19}}\;\text{GeV}
  \;=\; \frac{10^{3}}{1.22}\;\text{GeV}
  \;\approx\; 820\;\text{GeV}\,.
\]
\[
  \boxed{v \;\sim\; 820\;\text{GeV}\,.}
\]

\textbf{Comparison:} The measured value is
$v_{\text{measured}} = 246\;\text{GeV}$.  The estimate
$v \sim 820\;\text{GeV}$ differs by a factor of $\sim 3.3$,
which is within the $\mathcal{O}(1)$ uncertainty of the framework
(the coefficient $\sqrt{\xi/\lambda}$ absorbs this factor).

The parametric dependence $v \propto \Lambda_D^2$ is correct: if we
adjust $\Lambda_D$ downward by a factor of $\sim 1.8$, we get
$v = 246\;\text{GeV}$ exactly.  This is a parametric consistency
check, not a precision prediction.

\textit{(iii) Range of $\Lambda_D$ for $v \in [100, 1000]$~GeV.}

Inverting the seesaw relation:
\[
  \Lambda_D \;\sim\; \sqrt{v \cdot M_{\text{Pl}}}\,.
\]
For $v = 100\;\text{GeV}$:
\[
  \Lambda_D \;\sim\; \sqrt{100 \times 1.22 \times 10^{19}}
  \;=\; \sqrt{1.22 \times 10^{21}}
  \;\approx\; 3.5 \times 10^{10}\;\text{GeV}\,.
\]
For $v = 1000\;\text{GeV}$:
\[
  \Lambda_D \;\sim\; \sqrt{1000 \times 1.22 \times 10^{19}}
  \;=\; \sqrt{1.22 \times 10^{22}}
  \;\approx\; 1.1 \times 10^{11}\;\text{GeV}\,.
\]
\[
  \boxed{\Lambda_D \in
  [3.5 \times 10^{10},\; 1.1 \times 10^{11}]\;\text{GeV}
  \quad\text{for}\quad
  v \in [100,\; 1000]\;\text{GeV}\,.}
\]

This is a remarkably narrow range---less than one order of magnitude
in $\Lambda_D$ produces three orders of magnitude in $v$
(since $v \propto \Lambda_D^2$).

\textit{(iv) Dimensional consistency.}

\[
  \biggl[\frac{\Lambda_D^2}{M_{\text{Pl}}}\biggr]
  \;=\; \frac{[\text{energy}]^2}{[\text{energy}]}
  \;=\; [\text{energy}]\,.
  \quad\checkmark
\]

\bigskip

\textbf{(d) Summary and discussion.}

\textit{(i) Input/output table.}

\renewcommand{\arraystretch}{1.3}
\begin{center}
\begin{tabular}{@{} l l l @{}}
\toprule
\textbf{Quantity} & \textbf{Role} & \textbf{Value} \\
\midrule
$\Lambda_D$ & Input (free parameter) & $\sim 10^{11}\;\text{GeV}$ \\
$M_{\text{Pl}}$ & Input (measured) & $1.22 \times 10^{19}\;\text{GeV}$ \\
$\alphas(\Lambda_D)$ & Input (from duality relation/SM running)
  & $\sim 0.03$ \\
$v$ & Output (parametric consistency check) & $\sim 246\;\text{GeV}$ \\
$\Lambda_{\QCD}$ & Output (from RG running) & $\sim 200\;\text{MeV}$ \\
$\Lambda_S$ & Output (from spectrum) & $\sim \text{few TeV}$ \\
\bottomrule
\end{tabular}
\end{center}
\[
  \boxed{
  \text{Inputs: } \Lambda_D,\, M_{\text{Pl}},\, \alphas(\Lambda_D)\,.
  \qquad
  \text{Outputs: } v,\, \Lambda_{\QCD},\, \Lambda_S\,.}
\]

\textit{(ii) Energy hierarchy.}

The full hierarchy is:
\[
  \Lambda_{\QCD} \;\sim\; 200\;\text{MeV}
  \;\ll\; v \;\sim\; 246\;\text{GeV}
  \;\ll\; \Lambda_S \;\sim\; \text{few TeV}
  \;\ll\; \Lambda_D \;\sim\; 10^{11}\;\text{GeV}
  \;\ll\; M_{\text{Pl}} \;\sim\; 10^{19}\;\text{GeV}\,.
\]

\textbf{Inputs:}
\begin{itemize}
  \item $M_{\text{Pl}}$: fundamental gravitational scale (measured).
  \item $\Lambda_D$: duality scale (free parameter; constrained by
    consistency to $\sim 10^{10}$--$10^{11}\;\text{GeV}$).
  \item $\alphas(\Lambda_D)$: magnetic coupling at $\Lambda_D$ (from
    duality relation or evolved from measured $\alphas(M_Z)$).
\end{itemize}

\textbf{Outputs (derived):}
\begin{itemize}
  \item $v \sim \Lambda_D^2/M_{\text{Pl}}$: electroweak scale (seesaw
    relation).
  \item $\Lambda_{\QCD}$: from one-loop RG running below $\Lambda_D$.
  \item $\Lambda_S$: quasi-SUSY spectrum mass, of order
    $\Lambda_D^2/M_{\text{Pl}}$ (same seesaw as $v$, with different
    $\mathcal{O}(1)$ coefficients).
\end{itemize}

\textit{(iii) Prediction or consistency check?}

The relation $v \sim \Lambda_D^2/M_{\text{Pl}}$ is a
\textbf{parametric consistency check}, not a prediction:
\begin{itemize}
  \item \textbf{It is not a prediction} because the duality scale
    $\Lambda_D$ is a free parameter.  Given $\Lambda_D$, the framework
    produces $v$---but $\Lambda_D$ can be adjusted to match any desired
    $v$.
  \item \textbf{It is a consistency check} because the \emph{same}
    value of $\Lambda_D \sim 10^{11}\;\text{GeV}$ simultaneously
    gives: (a)~$v \sim 246\;\text{GeV}$ from the seesaw relation,
    (b)~$\Lambda_{\QCD} \sim 200\;\text{MeV}$ from RG running, and
    (c)~$\Lambda_S \sim \text{TeV}$ from the spectrum.  The fact that
    a single input $\Lambda_D$ produces three independent physical
    scales in the right parametric ballpark is non-trivial.
\end{itemize}
\[
  \boxed{\text{Parametric consistency check: one input }
  (\Lambda_D) \text{ gives three outputs }
  (v,\, \Lambda_{\QCD},\, \Lambda_S)
  \text{ in the correct range.}}
\]

\end{proof}


\bigskip

\begin{center}
\rule{0.6\textwidth}{0.4pt}

\textit{End of Exercise Set 6.  This exercise demonstrates the
scale-matching framework of the dualised Standard Model:
from a single input parameter $\Lambda_D \sim 10^{11}\;\text{GeV}$
(and the known Planck mass), the electroweak scale
$v \sim 246\;\text{GeV}$ emerges via a seesaw-like relation, and the
QCD confinement scale $\Lambda_{\QCD} \sim 200\;\text{MeV}$ follows
from one-loop RG running.  The numerical agreement---within
$\mathcal{O}(1)$ factors---is a parametric consistency check, not
a first-principles prediction.  The free parameter $\Lambda_D$ and the
$\mathcal{O}(1)$ coefficients from Planck-scale operators provide
the flexibility to match the observed scales exactly.}
\end{center}


%% ========================================================================
%%  BIBLIOGRAPHY
%% ========================================================================
\begin{thebibliography}{99}

\bibitem{CacciapagliaSannino2024}
  G.~Cacciapaglia and F.~Sannino,
  ``The Dualised Standard Model,''
  \texttt{arXiv:2407.17281} [hep-ph] (2024).

\bibitem{HillMachado2019}
  C.~T.~Hill, P.~A.~N.~Machado, A.~E.~Thomsen, and J.~Turner,
  ``Scalar democracy,''
  \textit{Phys.\ Rev.\ D} \textbf{100} (2019) 015015
  [\texttt{arXiv:1902.07214}].

\bibitem{Seiberg1995}
  N.~Seiberg,
  ``Electric--magnetic duality in supersymmetric non-Abelian gauge
  theories,''
  \textit{Nucl.\ Phys.} \textbf{B435} (1995) 129
  [\texttt{hep-th/9411149}].

\end{thebibliography}

\end{document}
